\section{Data and Measurement}\label{sec:data}

\subsection{Sample}
We build a firm-year panel from the XBRL financial data that companies file with the U.S.\ Securities and Exchange Commission \citep{sec2026xbrlapi}. The universe is every EDGAR registrant whose standard industrial classification (SIC) code is 7370 to 7374 (computer programming, data processing and related services). The EDGAR company index by SIC code lists inactive as well as active registrants, so firms later acquired or delisted remain in the universe and the sample does not condition on survival. One registrant is excluded from all data at the author's request; it is outside the estimation sample, and no figure reported here depends on it. For the 661 other registrants in this universe that ever report at least \$50 million of revenue, we retrieve every fact they filed in annual reports (forms 10-K, 20-F and 40-F, including amendments) for 35 US-GAAP elements.

Every item for a fiscal year is taken \emph{as first reported}: from the original annual filing for that year, the earliest-filed annual report that contains a revenue fact for the fiscal year. Growth is computed within that same filing, which reports the prior year alongside the current one. This matters. Later filings restate prior years, most often when a divested business is reclassified as discontinued operations, so mixing vintages produces spurious growth; for example, the restated continuing-operations revenue of AppLovin for 2023 is \AppRestated\% below its originally reported 2022 total, while growth as first reported was $+$\AppGrowth\%. Revenue is the largest revenue element reported for the year in that filing, which prevents the selection of a component. Items that the original filing does not report are taken from the earliest later filing that does; operating income, operating cash flow, capital expenditure and stock-based compensation come from the original filing in \VOrigLo{} to \VOrigHi\% of the sample firm-years. Fiscal years follow the Compustat convention: a year ending in January to May is assigned to the previous calendar year. We drop seven duplicate registrants (for example, a holding company and its operating subsidiary filing identical financial statements) and the overlapping years of three predecessor--successor pairs.

The sample is prepackaged software (SIC 7372), part of the industry to which the Rule of 40 is applied; the practitioner indexes cover software sold as a service, and SIC 7372 also includes game publishers and other firms outside that group, in fiscal years 2010 to 2025, with prior-year revenue of at least \$100 million and growth between $-90\%$ and $500\%$; margins outside $-300\%$ to $100\%$ are set to missing. In \VUsShare\% of firm-years the registrant has a U.S.\ business address; the rest are mostly foreign private issuers reporting under US GAAP, whose dollar growth includes currency movements.

\subsection{Sample rules}\label{sec:data-rules}
The review of the first version found firm-years that a frontier of software firms should not contain. We therefore define five rules, chosen before any estimate on the resulting sample was computed (Appendix~\ref{app:protocol}): four exclude firm-years and one flags them.
\begin{itemize}
\item \emph{R1, first reported.} A firm-year is excluded when its values come only from a filing whose own reporting period ends later, so that the year exists only as a comparative in a later report. This removes \VRone{} firm-years: \VRoneIpo{} years before a registrant's first annual report (comparatives in the first report after an initial public offering, spin-off or merger with a special-purpose acquisition company), \VRoneTen{} fiscal-2010 years before a registrant's first XBRL annual report and \VRoneLate{} years whose own report carries no revenue fact. Their median growth is \VRoneGrowth\%, against \VRoneGrowthOther\% for the other firm-years. The fiscal-2010 years of registrants that were already public were reported in real time, but the local data cannot separate them from pre-listing years, so all are excluded and added back in a sensitivity.
\item \emph{R2, principal business.} A firm whose principal business is not software is excluded. The rule reads the opening of the business section of each registrant's latest annual report (Item 1 of a 10-K, Item 4.B of a 20-F, the business section of the annual information form filed with a 40-F), which we retrieved for all \VRtwoDescribed{} registrants of the sample before the rules, and applies a keyword classifier that was fixed before any description was read. The classifier classes \VRtwoClassFirms{} firms as not software and excludes \VRtwoExclN{} firm-years: Bally Technologies and Inspired Entertainment (gambling) and Activision Blizzard, Changyou, Electronic Arts, GCL Global, PLAYSTUDIOS and Snail (games). It is a partial check: it misses game publishers and an AI-cloud provider that the review panels named, THQ, Take-Two, Roblox, Unity and CoreWeave, whose descriptions do not use its phrases. As a sensitivity, two readers, AI agents on different models that saw no outcome data, classified every description independently. They agree on \VRtwoReadAgree{} of the \VRtwoReadN{} firms (Cohen's $\kappa=\VRtwoKappa$) and class all \VRtwoClassFirms{} excluded firms as not software; both readers class \VRtwoReadBoth{} firms as not software, and at least one reader \VRtwoReadEither, including THQ, Take-Two, Roblox and CoreWeave and firms in IT services, hardware, telecommunications, payments, online education, advertising and e-commerce. Two sensitivities drop these firms as well (\VRtwoBothN{} and \VRtwoEitherN{} more firm-years), and a third drops all \VRtwoFirms{} named firms.
\item \emph{R3, captive revenue.} A firm-year is excluded when related parties provide at least \VCaptiveShare\% of revenue, measured by the XBRL related-party revenue elements and, where it was read, by the filing text. In the XBRL data \VRthreeXbrlN{} firm-years report related-party revenue, and none reaches the threshold (the largest share is \VRthreeMax\%, SINA in 2015). The filing text excludes \VRthree{} firm-year: Exzeo Group in 2025, whose annual report states that substantially all of its revenue came from customers affiliated with HCI. The XBRL element used most often was retired for recent years, so a captive filer since 2024 is caught only where its filing text was read.
\item \emph{R4, acquisitions.} The acquisition flag of Section~\ref{sec:data-acq} is extended in two ways: when the prior year-end goodwill is missing, the increase in goodwill is measured from the latest earlier year-end; and a cheap acquisition is flagged when cash paid or the goodwill increase is at least \VCheapShare\% of prior-year revenue and growth accelerates, relative to the year's median, by at least that amount, so that the acquisition could explain the acceleration at a price of at most one times revenue. On the earlier sample the flag rises from \VRfourOld{} to \VRfourNew{} firm-years. The rule was chosen among five variants after listing the firm-years each would flag and before any estimate on the new sample. A third criterion flags a year in which acquired businesses contributed at least 10\% of revenue since their acquisition dates (ASC 805-10-50-2(h)); it reaches \VRfourCN{} core firm-years, both already flagged, so it adds \VRfourCNew{} flags. Its coverage is thin, because the element is usually tagged on a business-combination axis that the SEC's frames omit.
\item \emph{R5, listed equity.} A filer without listed equity, such as a firm whose public debt keeps it filing after a take-private, is flagged rather than excluded. A 10-K firm-year is flagged when its cover page reports no public float outside the year of an initial public offering or spin-off: \VRfiveFlagN{} firm-years of \VRfiveFlagFirms{} firms, among them Medical Information Technology (Meditech), SunGard, Infor and Epicor. The flagged firm-years stay in the sample and a sensitivity drops them. Foreign filers are not tested, because their cover pages report shares outstanding rather than public float.
\end{itemize}
The resulting sample contains \VNEone{} firm-years for \VFirmsEone{} firms (\VNvTwo{} and \VFirmsvTwo{} before the rules): \VNPerTenFifteen{} in 2010--15, \VNPerSixteenNineteen{} in 2016--19, \VNPerTwentyTwentyTwo{} in 2020--22 and \VNPerTwentyThreeTwentyFive{} in 2023--25. Table~\ref{tab:robust} of the Supplementary Appendix reports every headline statistic with the excluded fiscal-2010 years added back, without the firm-years flagged under R5, without the firms that the readers of R2 class as not software and without the firms the review panels named.

\subsection{Growth and margin definitions}
Revenue growth is $g_{it}=S_{it}/S_{i,t-1}-1$. We study four margin definitions, each scaled by revenue:
\begin{enumerate}
\item GAAP operating margin, $m^{\mathrm{op}}$: operating income.
\item Adjusted operating margin, $m^{\mathrm{adj}}$: operating income plus SBC, close to the adjusted margins behind company-reported scores such as Palantir's \citep{palantir2026q2}.
\item Free-cash-flow margin, $m^{\mathrm{fcf}}$: operating cash flow (from continuing operations where reported) minus capital expenditure and capitalized software, which is read from the cash-flow statement where a filer tags it with an element of its own (Section~\ref{app:data} of the Supplementary Appendix). It is missing when capital expenditure is not reported (\VCapexMissing\% of firm-years) rather than set to zero.
\item Free-cash-flow margin after SBC, $m^{\mathrm{fcf-sbc}}$: free cash flow minus SBC, which treats equity compensation as the economic cost it is and corresponds to the free-cash-flow margin of the model.
\end{enumerate}
The Rule-of-40 score is $R=g+m$ under each definition, pairing growth over the past year with the same year's margin, the contemporaneous convention of reported scores (Section~\ref{sec:model}). US GAAP classifies interest paid and income taxes paid as operating cash flows (ASC 230-10-45-17(c) and (d)), so the two cash-flow scores are \emph{after} interest paid and cash taxes, while the GAAP operating margin is before both. Capital intensity is capital expenditure plus capitalized software over revenue. Following the model, we treat selling and marketing expense as the observable part of discretionary growth spend.

\subsection{Acquisitions and discontinued operations}\label{sec:data-acq}
Purchased growth is not the growth the model describes. We flag a firm-year as acquisition-affected when cash paid for acquisitions or the increase in goodwill over the fiscal year is at least 25\% of prior-year revenue, or when reported pro forma revenue exceeds actual revenue by at least 10\%; because acquired revenue arrives mostly in the following year, we flag the following year as well, matching years by fiscal year rather than by row. With the extensions of rule R4, the flag covers \VAcqShare\% of core firm-years and \VAcqShareFast\% of those with growth above 50\%. Acquisitions paid in stock and acquisitions whose only trace is the acquiree's revenue remain undetected. Discontinued operations are reported in \VDiscShare\% of core firm-years.

\subsection{Estimands}\label{sec:data-estimands}
A reported score and the frontier of the model answer different questions, so we define the estimand of every statistic.
\begin{itemize}
\item \emph{E1, observed scores (the baseline).} All \VNEone{} firm-years, with growth and margins as first reported.
\item \emph{E2, organic growth.} The model's growth is bought with expensed spend within one market, so acquired growth should be removed. Of the \VAffected{} acquisition-affected firm-years (\VAffAcqYear{} acquisition years and \VAffFollow{} following years), \VAffKept{} are kept with organic growth computed from the pro forma revenue that the original filings disclose. The other \VAffDropped{} (\VAffDroppedShare\%) cannot be adjusted and are dropped: in \VDropNoPf{} acquisition years the filing gives no pro forma revenue for the current year, in \VDropFollow{} following years the acquisition year's filing gives none usable, in \VDropAmbig{} the additive and cumulative readings of the disclosure differ by more than one point, and \VDropOther{} fail other checks of the disclosure. A further \VUnverified{} unflagged firm-years that grow by 100\% or more and whose filings report neither acquisition payments nor goodwill are left out of E2 and of the bounds, because their growth cannot be checked for acquisitions; \VUnaffected{} unaffected firm-years remain. E2 contains \VNEtwo{} firm-years for \VFirmsEtwo{} firms. For the \VAffKept{} adjusted firm-years the margin stays consolidated: it includes the acquired business, deal costs and, on the GAAP basis, amortization of acquired intangibles, so E2 pairs organic growth with a margin that is not organic; this lowers the GAAP margin of those firm-years and moves the cash margin in either direction.
\item \emph{Imputation bounds.} Because E2 drops most acquisition-affected years, it is not clear whether a difference between E1 and E2 reflects acquired growth or the years dropped. The bounds add those years back with organic growth imputed by removing acquired revenue equal to the purchase price over a price-to-revenue multiple $M$, half in the deal year and half in the next. At $M=4$ the median imputed acquired growth is \VImpFourAcq{} points, at $M=8$ \VImpEightAcq{} points (\VNImp{} firm-years in each). They bound acquired growth only for acquisitions priced at $M$ or more times revenue; for the cheap acquisitions that rule R4 flags they understate it.
\end{itemize}
E1 is the baseline. Every model-facing statistic, one that a result of Section~\ref{sec:model} speaks to, is reported under E1, E2 and both imputation bounds, and we state a model-facing claim only where its sign and its 90\% interval hold in all four (Table~\ref{tab:estimands_v3}, last column). Differences between E1 and E2 are shown side by side; Section~\ref{sec:empirics-robust} reports that none of them survives a familywise correction, and we make no claim that the estimands differ. The four views share most of their firm-years, so they are a sensitivity set, not four independent replications. Robustness checks also vary the flag threshold, the reading of ambiguous pro forma disclosures and the sample rules (Table~\ref{tab:robust}).

\subsection{Why the measurement basis matters}
The four definitions produce materially different scores for the same firm-year. In the sample, \VPassOp\% of firm-years clear 40 on the GAAP basis and \VPassCash\% after deducting SBC from free cash flow, against \VPassAdj\% on the adjusted basis and \VPassFcf\% on free cash flow before SBC. The bridge between the bases is large even at the median. In 2023--25 the median score is \VBridgeMedOp{} on the GAAP basis, \VBridgeMedAdj{} with SBC added back (median SBC is \VBridgeSbc\% of revenue), \VBridgeMedFcf{} on free cash flow before SBC and \VBridgeMedCash{} after SBC, and the shares clearing 40 are \VBridgePassOp\%, \VBridgePassAdj\%, \VBridgePassFcf\% and \VBridgePassCash\%. Adding back the interest and cash taxes that the cash score deducts raises its 95th percentile by \BridgeLiftLo{} to \BridgeLiftHi{} points in both 2016--19 and 2023--25 and leaves the decline between them in place (\BridgeChangePre{} points against \BridgeChangeBs{} on the same firm-years); adding back interest alone moves the 95th percentile by at most \BridgeIntMax{} points (Table~\ref{tab:bridge_v3} of the Supplementary Appendix). Removing interest income leaves a decline of similar size, although on the thin subsample that reports interest income the intervals are wide and the like-for-like interval includes zero (Section~\ref{app:data-additional} of the Supplementary Appendix). Interest and taxes paid come from the SEC's frames, which keep the latest filed value, so unlike the rest of the panel they are not first reported. Capital spending can change a score by more than 100 points. CoreWeave, which builds and operates AI computing capacity and which the review panels classified as compute infrastructure rather than software, grew \VCwGrowth\% in fiscal 2025 at a GAAP operating margin of about zero, for a GAAP score of \VCwOp; its free-cash-flow margin after SBC was \VCwFcfMargin\%, for a cash-basis score of \VCwCash. Company-reported scores such as Palantir's, which combine quarterly year-over-year growth with an adjusted quarterly margin, are not comparable with any of the four annual measures. We report the GAAP and free-cash-flow-after-SBC bases throughout and treat the latter as the economically relevant one.
